\section{Problem and available information}
The norm of the gradient need not vanish at a constrained minimum.
The full Frank-Wolfe gap accounts for the feasible directions and
does vanish there. We therefore use the gap in the curvature model,
derive a feasible descent step, and then use quasar-convexity to
convert descent into an objective-accuracy guarantee.

Let $Q\subset\mathbb R^n$ be nonempty, compact, and convex, with
$\operatorname{diam}Q\le D$, where $D>0$. Let $f$ be $C^2$ on a
neighborhood of $Q$, and fix a minimizer $x^*\in Q$.
All vector norms are Euclidean and matrix norms are spectral.
The available operations are evaluation of $f$, evaluation of its
gradient, and exact minimization of a linear function over $Q$.
At a current point $x$, compute
\begin{equation}\label{eq:gap}
 s\in\arg\min_{v\in Q}\langle\nabla f(x),v\rangle,\qquad
 d=s-x,\qquad
 G(x)=\langle\nabla f(x),x-s\rangle
     =\max_{v\in Q}\langle\nabla f(x),x-v\rangle.
\end{equation}
We assume
\begin{equation}\label{eq:gap_smoothness}
 \|\nabla^2f(x)\|\le L_0+L_1G(x)\quad(x\in Q),
 \qquad L_0,L_1\ge0.
\end{equation}
This condition depends on the feasible set. In particular,
$L_1D^2$ is dimensionless; the factor $D$ in the transport estimate
below cannot be dropped.

For convergence to the global minimum, we assume
$\gamma$-quasar-convexity about $x^*$:
\begin{equation}\label{eq:quasar}
 \gamma\bigl(f(x)-f(x^*)\bigr)
       \le\langle\nabla f(x),x-x^*\rangle\quad(x\in Q),
 \qquad 0<\gamma\le1.
\end{equation}
A convex objective satisfies this with $\gamma=1$. The condition
also admits nonconvex functions, as the example in Section~\ref{sec:examples}
shows. The step requires valid $L_0,L_1,D$; the objective-accuracy
test additionally uses $\gamma$. Neither $x^*$ nor $f(x^*)$ is an
input.

The constrained first-order condition gives $G(x^*)=0$, including
when $\nabla f(x^*)\ne0$. From \eqref{eq:gap} and \eqref{eq:quasar},
\begin{equation}\label{eq:certificate}
 0\le f(x)-f(x^*)\le \frac{G(x)}{\gamma}.
\end{equation}
Thus the gap bounds the objective error from above. Quasar-convexity
alone does not reverse this inequality.

\subsection{Relation to existing results and scope of the contribution}
The full gap and its use as an optimality certificate are standard~\cite{jaggi}.
The norm-dependent model and descent analysis of Vyguzov and
Stonyakin~\cite{vyguzov}, also published with Gasnikov~\cite{vyguzovjournal},
motivate our treatment of \eqref{eq:gap_smoothness}.
Mart\'{i}nez-Rubio~\cite{quasar} already studies constrained quasar-convex FW;
in particular, Theorem~9 gives $6LD^2/((k+1)\gamma^2)$ under ordinary
smoothness. Thus neither quasar-convex FW nor the factor $\gamma^{-2}$
is introduced here.

The accumulated certificate used below also has a known convex antecedent:
for $\gamma=1$ it is exactly the best primal-dual gap in
Wirth, Pe\~na, and Pokutta~\cite{wirth}, Section~2.2.
Their growth conditions relate the linear-model remainder to powers of $G$.
Section~\ref{sec:structure} gives an explicit connection to our assumption.
We use the certificate to analyze the proposed step for $0<\gamma\le1$;
we do not claim a new convex certificate.

Adaptive alternatives matter when comparing steps. Backtracking FW uses
objective tests to estimate curvature~\cite{pedregosa}; auto-conditioned FW
uses trial-point curvature without an inner search~\cite{acfw}.
Line-search-free conditional-gradient extensions also appear in
\cite{yagishita}. Adaptive Bregman steps handle a different, relative-smoothness
model~\cite{bregman}. These assumptions and guarantees are not identified
with \eqref{eq:gap_smoothness}.

\begin{table}[ht]
\centering\small
\caption{Position of the present analysis. Constants in different models
have different meanings; this is not a ranking of methods.}
\begin{tabular}{p{.23\linewidth}p{.34\linewidth}p{.34\linewidth}}
\toprule
Analysis & Curvature or growth information & Relevant distinction\\
\midrule
Norm-based FW~\cite{vyguzov} & $K_0+K_1\|\nabla f\|$ & Norm-dependent explicit short step\\
Quasar-convex FW~\cite{quasar} & Uniform $L$ & Known constrained quasar-convex rate\\
Open-loop FW~\cite{wirth} & Remainder controlled by $G^r$ & Known accumulated convex certificate\\
Adaptive FW~\cite{pedregosa,acfw} & Trial-point information & Does not require the supplied $L_0,L_1$\\
This analysis & $L_0+L_1G$ & Directional transport, explicit model minimizer, and certified stopping\\
\bottomrule
\end{tabular}
\end{table}

The contribution is the derivation and analysis of this particular curvature
model: a directional transport lemma, its logarithmic model minimizer,
separate $L_0/\varepsilon$ and $L_1\log(1/\varepsilon)$ contributions,
and verifiable examples. We also show that a conventional Armijo rule inherits
the same order of update complexity. The explicit rule saves trial function
evaluations when valid parameters are available; no universal numerical
advantage over adaptive FW is asserted.
